\documentclass[10pt,a4paper]{amsart}
\usepackage[margin=19mm]{geometry}
\usepackage{amsmath,amssymb,mathtools}
\usepackage[scr=rsfs,cal=euler]{mathalfa}
\usepackage{xfrac}
\usepackage[unicode,hidelinks,bookmarks=false]{hyperref}
\newcommand{\ie}{i.e.\ }
\newcommand{\cf}{cf.\ }
\newcommand{\resp}{respectively\ }
\newcommand{\Subsection}{\subsection}
\providecommand{\nobreakdash}{\nobreak\hbox{-}}
\setlength{\emergencystretch}{2em}
\multlinegap0pt
\usepackage{url,booktabs,tikz,setspace,fancyhdr,bm,mathrsfs}
\usepackage{enumitem}
\usepackage{cleveref}
\usepackage{bbm}
\usepackage{csquotes}
\usepackage{mathabx}
\usepackage{tikz}
\usepackage{float}
\counterwithin{figure}{section}
\newtheorem{theorem}{Theorem}[section]
\newtheorem{prop}[theorem]{Proposition}
\newtheorem{lemma}[theorem]{Lemma}
\newtheorem{cor}[theorem]{Corollary}
\newtheorem{conj}[theorem]{Conjecture}
\newtheorem{claim}[theorem]{Claim}
\newtheorem*{thm}{Theorem}
\newtheorem{fact}[theorem]{Fact}
\newtheorem{defn}[theorem]{Definition}
\newtheorem*{defn-non}{Definition}
\newtheorem{exam}[theorem]{Example}
\newtheorem{setup}[theorem]{Setup}
\newtheorem{exer}[theorem]{Exercise}
\newtheorem{ques}[theorem]{Question}
\newtheorem{rmk}[theorem]{Remark}
\newtheorem*{notation}{Notation}
\newenvironment{poc}{\begin{proof}[Proof of the claim]\renewcommand*{\qedsymbol}{$\blacksquare$}}{\end{proof}}
\newcommand{\BD}[1]{\boldsymbol{#1}}
\newcommand{\C}[1]{{\protect\mathcal{#1}}}
\newcommand{\B}[1]{{\bf #1}}
\newcommand{\I}[1]{{\mathbbm #1}}
\renewcommand{\O}[1]{\overline{#1}}
\newcommand{\M}[1]{\mathrm{#1}}
\newcommand{\ser}[1]{\mathsf{#1}}
\newcommand{\G}[1]{\mathfrak{#1}}
\newcommand{\V}[1]{\mathbold{#1}}
\newcommand{\cl}[1]{#1^{\cap}}
\newcommand{\Matrix}[2]{\left(\begin{array}{#1} #2\end{array}\right)}
\newcommand{\ceil}[1]{\lceil #1\rceil}
\newcommand{\floor}[1]{\lfloor #1\rfloor}
\newcommand{\innprod}[1]{\langle #1\rangle}
\newcommand{\me}{{\mathsf e}}
\newcommand{\aaa}{\boldsymbol{a}}
\newcommand{\bb}{\boldsymbol{b}}
\newcommand{\cc}{\boldsymbol{c}}
\newcommand{\dd}{\boldsymbol{d}}
\newcommand{\ee}{\boldsymbol{e}}
\newcommand{\vv}{\boldsymbol{v}}
\newcommand{\TT}{\boldsymbol{T}}
\newcommand{\RR}{\boldsymbol{R}}
\usepackage{todonotes}
\newcommand{\eps}{\varepsilon}
\newcommand{\ex}{\mathrm{ex}}
\newcommand{\cP}{\mathcal{P}}
\newcommand{\cQ}{\mathcal{Q}}
\newcommand{\cA}{\mathcal{A}}
\newcommand*{\bfrac}[2]{\genfrac{(}{)}{}{}{#1}{#2}}
\newcommand*{\abs}[1]{\lvert#1\rvert}
\newcommand*{\bigabs}[1]{\bigl|#1\bigr|}
\newcommand*{\Bigabs}[1]{\Bigl|#1\Bigr|}
\newcommand{\sY}{\mathscr{Y}}
\newcommand{\bT}{\boldsymbol{T}}
\newcommand{\bR}{\boldsymbol{R}}
\newcommand{\bv}{\boldsymbol{v}}
\newcommand{\bA}{\boldsymbol{A}}
\newcommand{\sD}{\mathscr{D}}
\newcommand{\sF}{\mathscr{F}}
\newcommand{\sH}{\mathscr{H}}
\newcommand{\shengtong}[1]{{\color{blue} #1}}
\newcommand{\cB}{\mathcal{B}}
\newcommand{\cD}{\mathcal{D}}
\newcommand{\cF}{\mathcal{F}}
\newcommand{\cG}{\mathcal{G}}
\newcommand{\cH}{\mathcal{H}}
\newcommand{\cS}{\mathcal{S}}
\newcommand{\cT}{\mathcal{T}}
\newcommand{\cX}{\mathcal{X}}
\newcommand{\cY}{\mathcal{Y}}
\newcommand{\cZ}{\mathcal{Z}}
\newcommand{\wj}[1]{{\color{magenta} #1}}
\newcommand{\VC}{\operatorname{VC}}
\newcommand{\Tr}{\operatorname{Tr}}
\newcommand{\AK}{\operatorname{AK}}
\newcommand{\W}{\operatorname{W}}
\begin{document}
\title[A disproof of the uniform witness conjecture]{A disproof of the uniform witness conjecture}
\author{Zixiang Xu}
\date{}
\maketitle
\noindent\emph{Source and licence.} A disproof of the uniform witness conjecture. Version arXiv:2606.24776v1 (CC BY 4.0). This material is distributed under the Creative Commons Attribution 4.0 International licence, http://creativecommons.org/licenses/by/4.0/, as recorded in the arXiv OAI licence metadata for this exact version and shown on the abstract page https://arxiv.org/abs/2606.24776v1. Full rights notice: licenses/arxiv-2606.24776-CC-BY-4.0.txt.\par\smallskip
\noindent\textit{Editorial scope.} Counted unit: the original \texttt{lemma} environment \texttt{lem:criterion} at source lines 198--203 together with its complete proof at lines 205--217; the delivered document keeps source lines 181--218 so that every referenced label and every citation resolves inside it. Native editable source is the paper's own concatenated TeX; the AMS wrapper is editorial formatting only. No publisher class, upstream script or unrelated artwork is redistributed.
\begin{theorem}\label{thm:main}
Let \(d\ge4\) and let \(s\) be an integer with
\(\left\lceil \frac{d+2}{2}\right\rceil\le s\le d-1.\) If \(n\ge2(d+1)\), then
\[
\W_{d,s}(n)\ge\binom{n-1}{d}+\binom{n-2(d+1-s)-2}{2s-d-2}.
\]
\end{theorem}
This covers every possible \(s\) when \(d\) is even, and every possible \(s\) except the lower boundary \(s=\frac{d+1}{2}\) when \(d\) is odd. Moreover, combining Theorem~\ref{thm:main} and the results in~\cite{2025CombProof,2026Chao,1961EKR} gives a rather strange picture: the conjecture is true for \(0\le s\le \frac{d}{2}\) and large \(n\), and it is also true for \(s=d\), but the middle parameters violate the conjectured bound.

\section{Proof of Theorem~\ref{thm:main}}

Let \(C\) and \(U\) be disjoint subsets of \([n]\), let \(\cP\subseteq 2^C\), and define
\[
\cF(\cP,U)=\bigg\{P\cup A:P\in\cP,\ A\in\binom{U}{d+1-|P|}\bigg\}.
\]
The set \(C\) is fixed and small. A set \(P\in\cP\) records the part of an edge inside \(C\), while \(A\subseteq U\) fills the remaining vertices.


\begin{lemma}\label{lem:criterion}
Let \(r=d+1-s\). Suppose that for every \(P\in\cP\), there is a set \(\tau(P)\subsetneq P\) such that \(|P|-|\tau(P)|\le r\), \(|\tau(P)|\le s\), and there is no \(Q\in\cP\) satisfying
\(
Q\cap P=\tau(P)\) and \(|Q|\le |\tau(P)|+r.\)
Then \(\cF(\cP,U)\) is an \(s\)-witness family.
\end{lemma}

\begin{proof}[Proof of Lemma~\ref{lem:criterion}]
For every \(P\in\cP\), we have \(|P|\le |\tau(P)|+r\le s+r=d+1\), so the family above is well-defined. Fix \(F=P\cup A\in\cF(\cP,U)\). Since \(|P|-|\tau(P)|\le r=d+1-s\), we get
\[
s-|\tau(P)|\le d+1-|P|=|A|.
\]
Choose \(X\subseteq A\) with \(|X|=s-|\tau(P)|\), and put \(B_F=\tau(P)\cup X\). Then \(B_F\subseteq F\) and \(|B_F|=s\). We claim that \(B_F\) is missing from the trace on \(F\). Indeed, suppose that \(F'=Q\cup A'\in\cF(\cP,U)\) satisfies \(F\cap F'=B_F\). Comparing the vertices in \(C\) and in \(U\) separately gives
\(Q\cap P=\tau(P)\) and \(A'\cap A=X.\)
Hence
\[
d+1-|Q|=|A'|\ge |X|=s-|\tau(P)|,
\]
which further implies that \(|Q|\le |\tau(P)|+r\), contradicting the assumption on \(\tau(P)\). Thus \(B_F\notin\Tr_{\cF}(F)\), and \(\cF(\cP,U)\) is an \(s\)-witness family.
\end{proof}


\begin{thebibliography}{99}
\bibitem[1961EKR]{1961EKR} P.~Erd\H{o}s, C.~Ko, and R.~Rado. \newblock Intersection theorems for systems of finite sets. \newblock {\em Quart. J. Math. Oxford Ser. (2)}, 12:313--320, 1961.
\bibitem[2025CombProof]{2025CombProof} T.-W. Chao, Z.~Xu, C.~H. Yip, and S.~Zhang. \newblock Uniform set systems with small {VC}-dimension. \newblock {\em Int. Math. Res. Not. IMRN}, 2025(17):Paper No. rnaf269, 20, 2025.
\bibitem[2026Chao]{2026Chao} T.-W. Chao, Z.~Xu, and D.~Zakharov. \newblock Uniform set systems with uniform witnesses. \newblock {\em arXiv preprint}, arXiv: 2602.17459, 2026.
\end{thebibliography}
\end{document}
